%=============================================================================!
% 3.8  THE ENERGY OF SEVERAL BODIES                                           !
%=============================================================================!
\section{The energy of several bodies}
\label{sec:en-bodies}

\Cref{sec:ne-bodies} showed that the forces two bodies exert on each
other cannot change their total momentum. This section asks the same
question of energy, using the two carts of that section, and finds a
different answer: forces between bodies can change their total kinetic
energy, and they conserve the total energy only when they store what they
give in a potential energy.

\subsection*{The pair of forces does work}

Two carts, of masses $m_1 = \SI{1}{\kilogram}$ and $m_2 =
\SI{3}{\kilogram}$, at positions $x_1 < x_2$ on a level, smooth track,
are pushed apart by a light spring of stiffness $k =
\SI{200}{\newton\per\metre}$ compressed between them
(\cref{fig:ne-carts-sketch}). Let $\xi$ be the compression of the spring,
how much shorter it is than its natural length. The spring touches both
carts, so as the gap $x_2 - x_1$ between them grows, $\xi$ shrinks by the
same amount:
\begin{equation*}
   \frac{\dd \xi}{\dd t} = -\frac{\dd}{\dd t}(x_2 - x_1) = -(v_2 - v_1) .
\end{equation*}
By \cref{sec:ne-bodies} the spring pushes the two carts with forces of
equal size in opposite directions, and by Hooke's law
(\cref{sec:ne-spring}) a spring compressed by $\xi$ pushes with a force of
size $k\xi$, so it pushes the first cart with $-k\xi$ and the second with
$+k\xi$. The weights and normal forces do no work, so by \cref{sec:en-varying} the kinetic energies change at the rates
\begin{equation*}
   \frac{\dd K_1}{\dd t} = -k\xi\,v_1 , \qquad
   \frac{\dd K_2}{\dd t} = +k\xi\,v_2 .
\end{equation*}
Adding them,
\begin{align*}
   \frac{\dd}{\dd t}\bigl(K_1 + K_2\bigr) &= k\xi\,(v_2 - v_1)\\
   &= -k\xi\,\frac{\dd \xi}{\dd t}
      && \text{using the rate of change of $\xi$}\\
   &= -\frac{\dd}{\dd t}\bigl(\tfrac12 k\xi^2\bigr)
      && \text{since $\tfrac{\dd}{\dd t}\bigl(\tfrac12 k\xi^2\bigr)
         = k\xi\,\tfrac{\dd \xi}{\dd t}$ by the chain rule.}
\end{align*}
Moving the last term to the left, the sum
\begin{equation*}
   E = \tfrac12 m_1 v_1^2 + \tfrac12 m_2 v_2^2 + \tfrac12 k \xi^2
\end{equation*}
does not change: the total energy of the two carts and the spring is
conserved, and $\frac12 k\xi^2$ is the potential energy of the spring, as in
\cref{sec:en-potential}. It depends only on how far apart the carts are.

The contrast with momentum lies in the velocities. The two forces are
equal and opposite, so they cancel in the total momentum, whose rate of
change is their sum. Their powers are $-k\xi\,v_1$ and $+k\xi\,v_2$, which
cancel only if $v_1 = v_2$. When the carts move apart, the first moves
backwards, along its push, and the second forwards, along its push, so
both forces do positive work and the total kinetic energy grows. A pair
of equal and opposite forces does work whenever the distance between the
bodies changes. This is how the carts, at rest at first, come to move.

\subsection*{The speeds of the carts}

The spring starts compressed by \SI{0.1}{\metre}, so it stores $\frac12
\times200\times0.1^2 = \SI{1}{\joule}$, and the carts start at rest. Once
the spring has reached its natural length and fallen away, $\xi = 0$ and
all of that energy is kinetic,
\begin{equation*}
   \tfrac12 m_1 v_1^2 + \tfrac12 m_2 v_2^2 = \SI{1}{\joule} .
\end{equation*}
One equation cannot fix two speeds, but momentum supplies the second. The
total momentum stays zero (\cref{sec:ne-bodies}), so $m_1 v_1 + m_2 v_2
= 0$, which gives $v_1 = -(m_2/m_1)\,v_2 = -3v_2$. Substituting,
\begin{equation*}
   \tfrac12\times1\times(-3v_2)^2 + \tfrac12\times3\times v_2^2
   = \tfrac92 v_2^2 + \tfrac32 v_2^2 = 6\,v_2^2 = 1 ,
\end{equation*}
in joules, so $v_2^2 = \frac16$ and
\begin{equation*}
   v_2 = +\sqrt{\tfrac16} = +\SI{0.408}{\metre\per\second} , \qquad
   v_1 = -3\sqrt{\tfrac16} = -\SI{1.225}{\metre\per\second} ,
\end{equation*}
taking the positive root for $v_2$ because the spring pushes the second
cart forwards. These are the speeds that a computer solution of the
equations of motion gave in \cref{sec:ne-bodies}, found now with two
lines of algebra. \Cref{fig:en-carts} shows the energy passing from the
spring to the carts while the total stays at \SI{1}{\joule}.

\begin{figure}[htb]
\centering
\includegraphics{ch03_energy/carts}
\caption{The energies of the carts of \cref{fig:ne-carts}: the energy
stored in the spring, $\frac12 k\xi^2$, the kinetic energies of the two
carts, and the total, which stays at \SI{1}{\joule}. The spring falls away
at the dotted line, after \SI{0.096}{\second}.}
\label{fig:en-carts}
\end{figure}

The lighter cart ends with \SI{0.75}{\joule} and the heavier with
\SI{0.25}{\joule}, although their momenta are equal and opposite. Writing
the kinetic energy in terms of the momentum $p = mv$ shows why,
\begin{equation*}
   K = \tfrac12 m v^2 = \frac{(mv)^2}{2m} = \frac{p^2}{2m} :
\end{equation*}
for momenta of equal size the kinetic energies are in the inverse ratio of
the masses, three to one here.

The same argument extends to any number of bodies that push on one
another in third-law pairs, provided that the two forces of each pair
act along the line joining its two bodies, with a size that depends only
on the distance between them: each pair then has a potential energy, and the kinetic energy plus the sum of the pair energies
is conserved. \Cref{sec:en-atoms} proves this for atoms, in three
dimensions.

\notebook{03}{change the masses of the carts and the energy they share}

\begin{selfcheck}
If both carts had a mass of \SI{2}{\kilogram}, with the same spring and
compression, at what speeds would they separate?
\end{selfcheck}
